\begin{document}

\begin{center}
{\Large\textbf{Объём как доля целого}}\\[2mm]
{\small ЕГЭ, стереометрия, задание №3 \quad | \quad Фамилия\ \answerline{6cm}}
\end{center}

\ifteacher
\begin{center}\fbox{\textbf{Учительская версия:} ответы и краткие опоры выделены курсивом.}\end{center}
\fi

\textbf{Цель.} Не искать координаты: сначала вырази искомое тело как долю
исходного, разность или сумму знакомых призм и пирамид. Рисунок является
схемой: измерять на нём ничего не нужно.

\section*{1. Формула объёма пирамиды}

\stereotask{1}{Найдите объём правильной треугольной пирамиды, стороны
основания которой равны \(12\), а высота равна \(6\sqrt3\).}{\csname figStereoV28\endcsname}
{Площадь равностороннего основания \(36\sqrt3\). Ответ:
\(36\sqrt3\cdot6\sqrt3/3=216\).}

\stereotask{2}{Сторона основания правильной шестиугольной пирамиды равна
\(3\), боковое ребро равно \(6\). Найдите объём пирамиды.}{\csname figStereoV04\endcsname}
{Радиус описанной окружности основания \(3\), высота \(\sqrt{6^2-3^2}=3\sqrt3\).
Площадь основания равна \(27\sqrt3/2\), ответ \(81/2=40{,}5\).}

\section*{2. Одна и та же высота: объём как доля}

\stereotask{3}{От треугольной призмы объёмом \(120\) отсечена треугольная
пирамида плоскостью, проходящей через сторону одного основания и
противоположную вершину другого основания. Найдите объём оставшейся части.}
{\csname figStereoV07\endcsname}
{Отсечённая пирамида составляет треть объёма призмы: \(40\). Ответ: \(80\).}

\stereotask{4}{Объём треугольной пирамиды равен \(14\). Плоскость проходит
через сторону основания и пересекает противоположное боковое ребро в точке,
делящей его в отношении \(2:5\), считая от вершины. Найдите больший из
получившихся объёмов.}{\csname figStereoV08\endcsname}
{Объёмы относятся как высоты: \(2:5\). Большая часть равна
\(14\cdot5/7=10\).}

\stereotask{5}{Объём параллелепипеда \(ABCDA_1B_1C_1D_1\) равен \(60\).
Найдите объём треугольной пирамиды \(ACB_1D_1\).}{\csname figStereoV14\endcsname}
{Пирамида составляет треть объёма параллелепипеда. Ответ: \(20\).}

\stereotask{6}{В прямоугольном параллелепипеде \(ABCDA_1B_1C_1D_1\) известно,
что \(AB=9\), \(BC=8\), \(AA_1=6\). Найдите объём многогранника с вершинами
\(A\), \(B\), \(C\), \(B_1\).}{\csname figStereoV16\endcsname}
{Площадь основания \(\triangle ABC\) равна \(9\cdot8/2=36\), высота \(6\).
Ответ: \(36\cdot6/3=72\).}

\section*{3. Сумма и разность простых тел}

\stereotask{7}{В прямоугольном параллелепипеде \(ABCDA_1B_1C_1D_1\) известно,
что \(AB=9\), \(BC=6\), \(AA_1=5\). Найдите объём многогранника с вершинами
\(A\), \(B\), \(C\), \(D\), \(A_1\), \(B_1\).}{\csname figStereoV15\endcsname}
{Это половина параллелепипеда: \(9\cdot6\cdot5/2=135\).}

\stereotask{8}{Найдите объём многогранника с вершинами \(A\), \(B\), \(C\),
\(D\), \(B_1\) прямоугольного параллелепипеда \(ABCDA_1B_1C_1D_1\), если
\(AB=9\), \(BC=3\), \(BB_1=8\).}{\csname figStereoV33\endcsname}
{Многогранник --- пирамида с основанием \(ABCD\): \(9\cdot3\cdot8/3=72\).}

\section*{4. Сечения через середины}

\stereotask{9}{Объём треугольной призмы, отсекаемой от куба плоскостью через
середины двух рёбер, выходящих из одной вершины, и параллельной третьему
ребру, равен \(25\). Найдите объём куба.}{\csname figStereoV25\endcsname}
{Отсечённая призма составляет \(1/8\) объёма куба. Ответ: \(200\).}

\stereotask{10}{Объём треугольной призмы, отсекаемой от куба плоскостью через
середины двух рёбер, выходящих из одной вершины, и параллельной третьему
ребру, равен \(11\). Найдите объём куба.}{\csname figStereoV26\endcsname}
{Отсечённая призма составляет \(1/8\) объёма куба. Ответ: \(88\).}

\stereotask{11}{От треугольной пирамиды объёмом \(42\) отсечена треугольная
пирамида плоскостью, проходящей через вершину исходной пирамиды и среднюю
линию основания. Найдите объём отсечённой пирамиды.}{\csname figStereoV27\endcsname}
{Площадь основания отсечённой пирамиды составляет \(1/4\), высота та же.
Ответ: \(42/4=10{,}5\).}

\ifteacher
\section*{Короткая проверка урока}
Ученик готов к следующему блоку, если до вычислений записывает долю искомого
объёма и может объяснить её через равные высоты, площади оснований или
разбиение исходного тела.
\fi

\end{document}
